% !TeX spellcheck = en_US
\documentclass[french]{yLectureNote}

\title{Électromagnétisme 3}
\subtitle{Physique}
\author{Paulhenry Saux}
\date{\today}
\yLanguage{Français}
%jesse.groenen@cemes.fr
\professor{M.Brut}
\usepackage{graphicx}%----pour mettre des images
\usepackage[utf8]{inputenc}%---encodage
\usepackage{geometry}%---pour modifier les tailles et mettre a4paper
%\usepackage{awesomebox}%---pour les boites d'exercices, de pbq et de croquis ---d\'esactiv\'e pour les TP de PC
\usepackage{tikz}%---pour deiffner + d\'ependance de chemfig
% \usepackage{tabularx}%---pour dimensionner automatiquement les tableaux avec variable X
\usepackage{awesomebox}%---Pour les boites info, danger et autres
\usepackage{menukeys}%---Pour deiffner les touches de Calculatrice
\usepackage{fancyhdr}%---pour les en-t\^ete personnalis\'ees
\usepackage{blindtext}%---pour les liens
\usepackage{hyperref}%---pour les liens (\`a mettre en dernier)
\usepackage{caption}%---pour la francisation de la l\'egende table vers Tableau
\usepackage{pifont}
\usepackage{array}%---pour les tableaux
\usepackage{lipsum}
\usepackage{yFlatTable}
\usepackage{multicol}
\usepackage{tabularx}
\usepackage{booktabs}
\newcommand{\Lim}[1]{\lim\limits_{\substack{#1}}\:}
\renewcommand{\vec}{\overrightarrow}
\newcommand{\N}[0]{\mathbb{N}}
\newcommand{\dd}{\mathrm{d}}
\newcommand{\norm}[1]{||\vec{#1}||}
\newcommand{\fo}{\psi(\vec{r},t)}
\newcommand{\foe}{\psi(\vec{r},t)\*}
\newcommand{\HH}{\hat{H}}
\newcommand{\hb}{\hbar}
\newcommand{\lap}{\nabla^2}
\newcommand{\lapcc}{\frac{\partial^2 }{\partial x^2}+\frac{\partial^2 }{\partial y^2}+\frac{\partial^2 }{\partial z^2}}
\newcommand{\E}{\vec{E}(M)}
\newcommand{\B}{\vec{B}(M)}
\newcommand{\V}{V(M)}
\newcommand{\er}{\vec{e_{\rho}}}
\newcommand{\ep}{\vec{e_{\varphi}}}
\newcommand{\pip}{\pi^+}
\newcommand{\pim}{\pi^-}
\newcommand{\mv}{\mathcal{V}}
\DeclareMathOperator{\Div}{Div}
\newcommand{\rot}{\vec{\mathrm{rot}}}
\newcommand{\grad}{\vec{\mathrm{grad}}}
%\setcounter{chapter}{3}
\begin{document}
%\chapter{Puissance et énergie du champ Em}
\setcounter{chapter}{5}
\chapter{Aspects microscopique}
\section{Moment magnétiques atomiques}
\subsection{Moment magnétique orbital}
\begin{definition}[Moment magnétique d'un électron autour d'un noyau]
    Moment d'un électron autour d'un noyau : $\vec{L} = \vec{r}\wedge (m_e \vec{v})$

    Moment d'un courant circulaire : $\vec{\mu} = iS\vec{n} = i\pi r^2\vec{n}$

    Le moment magnétique total : $\vec{\mu} = -\frac{e}{2m_e}\vec{L} = \gamma_e \vec{L}$
\end{definition}
\begin{definition}[Magnéton de Bohr]
    $\mu_{l,z} =\gamma_e L_z =  -m_l \mu_B$ avec $\mu_B$ le magnéton de Bohr.
\end{definition}
\subsection{Moment magnétique de Spin}
On ajoute le moment cinétique intrinsèque de l'électron : Spin S : $S_z = m_s \hbar$
\subsection{Moment total}
$\vec{j} = \vec{L} + \vec{S}$ et $\norm{J}^2 = (j(j+1))\hbar^2$
\section{Diamagnétisme}
On étudie le phénomène d'induction à l'échelle microscopique en présence d'un champ magnétique appliqué : $\vec{B_a} = B_a \vec{e_z}$
\begin{theorem}{Orientation d'aimantation pour le diamagnétisme}
    L'orientation de l'aimantation est opposée à celle de $\vec{B_a}$, conformément à la Loi de Lenz. Son effet est très faible
\end{theorem}

Le diamagnétisme est extrêmement faible : $\chi_m \simeq 10^{-6}$

Dans ce cas, $\vec{B_{local}}\simeq \vec{B_a}$
% \section{Paramagnétisme}
% Ici, $\chi_m \simeq 10^{-3}$


% Il faut un moment cinétique total $\vec{j}\neq \vec{0}$

% Le Paramagnétisme dépend de la température
% \subsection{Théorie de Brillouin}
% \subsubsection{Théorie}
% %TODO Explication très simple
% \subsubsection{Cas simple : 1 é célibataire}
% On a $\vec{M} = n_v \mu_b \th(X)\vec{e_z}$

% Quand $X>> 1, \th(X)\to 1$ : Tous les moments sont alignés sur $\vec{B_a}$. C'est le cas pour un régime en basse température ou un champ très fort.

% Quand X est faible, on a $\th(X) \simeq X$ : on a une expression linéaire.

% Si $n_v$ est trop grand, l'approximation des moments indépendants cesse d'être valide



\section{Paramagnétisme}

Le paramagnétisme désigne le comportement d'un milieu matériel qui n'a pas d'aimantation spontanée, mais qui acquiert sous l'effet d'un champ magnétique extérieur $\vec{B_a}$ une aimantation alignée dans le même sens que ce champ.

\checkInfo{Ordre de grandeur}{La susceptibilité magnétique est positive et faible : $\chi_m \simeq 10^{-3}$ à $10^{-5}$.}

\criticalInfo{Condition atomique}{Pour observer du paramagnétisme, il faut que les atomes ou ions possèdent un moment cinétique total non nul ($\vec{J} \neq \vec{0}$), ce qui implique la présence d'électrons célibataires (couches électroniques incomplètes).}

\subsection{Théorie de Brillouin}

\subsubsection{Théorie}

L'origine du paramagnétisme est quantique. Lorsqu'un atome possède un moment magnétique permanent $\vec{\mu}$, son interaction avec un champ magnétique extérieur $\vec{B_a}$ est décrite par l'énergie potentielle dipolaire :
$$E_p = -\vec{\mu} \cdot \vec{B_a}$$
\subsubsection{Cas simple : un électron célibataire}

\begin{proposition}
L'aimantation macroscopique $\vec{M}$ d'un système de $n_v$ moments indépendants par unité de volume s'exprime par :
$$\vec{M} = n_v \mu_B \tanh(X)\vec{e_z}$$
où $\mu_B$ est le magnéton de Bohr, $\vec{e_z}$ est la direction du champ appliqué $\vec{B_a}$, et $X$ est le paramètre sans dimension :
$$X = \frac{\mu_B B_a}{k_B T}$$
\end{proposition}


\subsubsection{Gas général}
\begin{definition}
    Aimantation volumique : $\vec{M} = n_v g\mu_B B_j(X)\vec{e_z}$

\end{definition}
\section{Ferromagnétisme}
\begin{definition}
    Cela résulte de l'interaction entre atomes du fait des spins des électrons des couches périphériques incomplètes associées à des
sites voisins.

On peut avoir une aimantation même en l'absence de champ appliqué. On qualifiera cette aimantation de spontannée.
\end{definition}

% \subsection{Courbe de première aimantation et cycle Histérésis}
% M ne prend pas les mêmes valeurs durant l'évolution de H. C'est un effet d'histérésis.

% On écrit : $\mu_r = \frac{B}{\mu_0H}$
% \subsection{Interpération microscopique}
% \begin{definition}
%     Modèle de Weiss : $\vec{B} = \vec{B_a}+\lambda \vec{M}$
% \end{definition}

\subsection{Courbe de première aimantation et cycle d'hystérésis}

Pour les matériaux ferromagnétiques, la relation entre l'aimantation $\vec{M}$ (ou l'induction $\vec{B}$) et le champ excitateur $\vec{H}$ n'est plus linéaire et dépend de l'histoire magnétique du matériau.

\begin{definition}[Courbe de première aimantation]
    C'est la courbe obtenue en augmentant progressivement le champ $H$ à partir d'un état initial totalement désaimanté ($H=0, M=0$). Elle présente une zone de forte croissance avant de stagner à la saturation $M_{\text{sat}}$.
\end{definition}

Lorsqu'on fait osciller le champ $H$ entre deux valeurs extrêmes $-H_{\max}$ et $+H_{\max}$, $M$ ne suit pas le même chemin à l'aller et au retour : c'est le phénomène d'hystérésis.

\warningInfo{Caractéristiques du cycle}{On définit deux grandeurs fondamentales sur un cycle d'hystérésis :
\begin{itemize}
    \item \textbf{L'aimantation rémanente $M_r$ (ou $B_r$) :} L'aimantation subsistant lorsque le champ appliqué est annulé ($H = 0$).
    \item \textbf{Le champ coercitif $H_c$ :} Le champ inverse qu'il faut appliquer pour annuler l'aimantation ($M = 0$).
\end{itemize}}

\subsection{Interprétation microscopique}
\begin{definition}
    \textbf{Modèle du champ moléculaire de Weiss :} On modélise ces interactions quantiques complexes (interactions d'échange) en postulant que chaque moment magnétique subit, en plus du champ extérieur $\vec{B_a}$, un champ effectif proportionnel à l'aimantation du milieu, appelé champ moléculaire :
    $$\vec{B}_{\text{eff}} = \vec{B_a} + \lambda \vec{M}$$
    où $\lambda$ est la constante de Weiss ($\lambda > 0$).
\end{definition}
\end{document}
